% Elaborado por Fabian Gomez
% Version 1.0

\section{Introduction}
 Este documento constituye la Especificación de Requisitos del Sistema (\textit{System Requirements Specification}, SyRS/SRS) del Rover Olympus y establece un baseline único y trazable de requisitos, arquitectura e información de Verificación y Validación (V\&V). El contexto del sistema, su propósito como plataforma de validación del software de navegación autónoma ELANav y el rol del UC-01 (Exploración de superficie) como caso núcleo se presentan en §\ref{sec:system_context}; la definición completa de casos de uso top-level se encuentra en §\ref{subsec:use_cases_top_level} y el detalle del UC-01 se desarrolla en §\ref{sec:uc01}.

\textbf{Linaje técnico y delta respecto al predecesor POAS.} El Rover Olympus se concibe como evolución directa del rover POAS (\textit{Prototipo de Observación y Análisis Superficial}), desarrollado previamente en el SETEC Lab/TEC por el grupo TECSpace \cite{tecspace2021poas}. POAS validó mecánicamente una suspensión \textit{rocker-bogie} en terreno volcánico local, pero su arquitectura electrónica era predominantemente reactiva/teleoperada, con percepción y cómputo embarcado limitados. Olympus capitaliza la arquitectura mecánica heredada de POAS (suspensión rocker-bogie, chasis tubular híbrido PVC/3D) e introduce el \textbf{delta} necesario para soportar la validación de navegación autónoma del proyecto ELANav: (i) electrónica de potencia con drivers mixtos L298N+BTS7960 y telemetría INA3221 por bancos, (ii) capa de percepción activa con cámara OV5647 + ToF (VL53L0X) + LiDAR (TF02) + ultrasonidos HC-SR04, (iii) cómputo embarcado de doble nivel (Arduino Mega LLC + Raspberry Pi~5 HLC con stack \texttt{olympus\_hlc} sobre Yocto), y (iv) protocolo CSP sobre WiFi para enlace rover--GCS. Esta extensión transforma a POAS de banco de pruebas mecánico a plataforma de validación de software autónomo, satisfaciendo el Objetivo General del TFG.

\textbf{Origen de los casos de uso.} Los casos de uso top-level (UC-01\dots UC-05) se derivan de escenarios representativos de una misión lunar de exploración de superficie y se reducen a su validación equivalente en sitio análogo terrestre conforme al alcance académico del proyecto \cite{ieee_29148, nasa_se_handbook}.


\subsection{Purpose}

El propósito de este SyRS es definir de manera completa, verificable y trazable los requisitos del Rover Olympus como plataforma de validación para ELANav, orientada a demostrar capacidades de navegación autónoma sobre terreno en condiciones análogas, conforme al entorno operacional descrito en §\ref{subsec:operational_environment} y a las restricciones operacionales de sitio análogo en §\ref{subsec:operational_constraints}. \cite{ieee_29148}

Este documento sirve como referencia técnica principal para:
\begin{itemize}
    \item \textbf{Guiar el diseño e integración del sistema:} Hardware Commercial Off-The-Shelf (COTS) y software embebido.
    \item \textbf{Establecer criterios de desempeño y calidad:} Criterios medibles para UC‑01.
    \item \textbf{Planificar y ejecutar V\&V:} Soporte para campañas de prueba representativas. \cite{nasa_se_handbook, ieee_29148}
\end{itemize}

Adicionalmente, este SyRS formaliza la coherencia entre ConOps, Casos de Uso, arquitectura lógica/física (7 subsistemas) y la matriz de cumplimiento V\&V. \cite{nasa_se_handbook}

\subsection{Scope}

\textbf{Alcance del sistema:} El Rover Olympus es un rover terrestre de exploración de carácter académico, integrado principalmente con hardware COTS y diseñado como plataforma de validación de ELANav. La descripción del sistema, su frontera y entidades externas se detallan en §\ref{subsec:system_description}–§\ref{subsec:system_boundary}. \cite{ieee_29148}

\textbf{Alcance funcional:} Este SyRS especifica los requisitos necesarios para soportar los casos de uso top-level del proyecto, manteniendo UC-01 como caso núcleo de operación y validación. La lista y descripción de los casos de uso UC-01…UC-05 se presenta en §\ref{subsec:use_cases_top_level}, y el detalle completo de UC-01 se desarrolla en §\ref{sec:uc01}. \cite{nasa_se_handbook}

\textbf{Alcance de validación:} La validación se ejecuta en sitio análogo terrestre; el entorno operacional se caracteriza en §\ref{subsec:operational_environment} y las restricciones operacionales específicas se enumeran en §\ref{subsec:operational_constraints}. Los criterios y métricas de aceptación/verificación se consolidan en las secciones de requisitos y V\&V correspondientes (§\ref{sec:requirements}–§\ref{sec:vv}). \cite{nasa_se_handbook}

\subsection{Intended Audience and Use}

Este SyRS está dirigido a:
\begin{enumerate}[label=\alph*)]
    \item Investigadores y desarrolladores del SETEC Lab/TEC responsables de ELANav y del rover.
    \item Ingenieros de integración y pruebas que ejecutan campañas en sitio análogo.
    \item Evaluadores académicos que revisan cumplimiento y evidencia de V\&V. \cite{ieee_29148}
\end{enumerate}

El uso previsto es:
\begin{enumerate}
    \item Especificar requisitos verificables del sistema centrados en UC‑01.
    \item Servir como baseline técnico para decisiones de arquitectura (Model-Based Systems Engineering, MBSE/SysML) e interfaces.
    \item Soportar planificación y ejecución de V\&V mediante criterios cuantitativos y métodos T/A/D/I. \cite{nasa_se_handbook, segoldmine_ppi}
\end{enumerate}

\subsection{Document Overview}

El documento incluye:
\begin{itemize}
    \item Referencias normativas, términos y acrónimos.
    \item Contexto del sistema y entorno operacional.
    \item Stakeholders, actores y resumen ConOps.
    \item Interfaces y organización de requisitos.
    \item Requisitos del sistema (RF/RNF/CON) agrupados por categorías.
    \item Referencias de arquitectura y modelos (MBSE/SysML).
    \item Estrategia de V\&V, criterios de aceptación y trazabilidad. \cite{ieee_29148}
\end{itemize}

La sección 11 contiene la ``definición maestra'' de los requisitos RF/RNF/CON; otras secciones solo referencian IDs cuando sea necesario para contexto o trazabilidad. \cite{nasa_se_handbook, segoldmine_ppi}

\subsection{Conformance and Tailoring Statement}
Este SyRS adopta la estructura de un SRS/SyRS conforme a ISO/IEC/IEEE 29148:2018 y se alinea con prácticas NASA de ingeniería de sistemas para proyectos de exploración robótica. Se aplica un proceso de \textit{tailoring} (adaptación) justificado por el contexto académico y los objetivos de validación de ELANav: \cite{ieee_29148, nasa_se_handbook}

\begin{itemize}
    \item \textbf{Fusión SyRS/SRS:} Se consolidan los requisitos de sistema (SyRS) y software (SRS) en un único baseline técnico para garantizar la trazabilidad directa y agilidad en una plataforma de validación rápida.
    \item \textbf{Priorización de UC-01:} Se focaliza el desarrollo detallado de requisitos en el caso de uso núcleo (Exploración), estableciendo un baseline verificable para la misión primaria, mientras que UC-02 a UC-05 se mantienen en nivel conceptual.
    \item \textbf{Gestión centralizada de TBDs:} Se utiliza un registro de parámetros pendientes (Anexo B) con un plan de resolución operativa, permitiendo la evolución del documento conforme se caracterizan los componentes COTS en campo.
\end{itemize} 
